\chapter*{前書き}


これは平面幾何学を網羅的に述べようとするものではありません。
作者の私的なまとめです。

一階述語論理の幾何学の公理系としてはタルスキによるものがありますが、公理系がいわば極めて圧縮されているので、私としては非常に扱いにくいです。

こちらのものは、公理の数が多く、しかも公理 \ref{261002134501} が長たらしいですが、公理と題した部分は「単に、義務教育の数学で経験的に身につける常識を形式化したもの」と割り切った部分であり、主眼は公理よりも公準にあります。

ただし、公理 \ref{261002134501} については補足しておくことがあります。
この公理は実は、「正の数には必ず平方根が存在する」と「奇数次の一変数方程式には必ず解が存在する」という二つの公理に置き換えることができます。
そうしない理由は、等価であることの証明が厄介なのと、「実数体の公理系をそのまま一階述語論理で表現し直す」という基本の考えが見えにくくなるからです。
